\section{Классы языков}
	\tikzsetfigurename{FLT_LangClasses_}

\begin{marginfigure}
    \begin{center}
        \resizebox{\marginparwidth}{!}{
                        \begin{tikzpicture}
                \node at (3,0) (reg) {Регулярные языки (Тип 3)};
                \node at (3,1) (cfl) {Контекстно-свободные языки (Тип 2)};
                \node at (3,2) (cdl) {Контекстно-зависимые языки (Тип 1)};
                \node at (3,3) (arb) {Неограниченные языки (Тип 0)};

                \node[r_state,fit=(reg)] (a) {};
                \node[r_state,fit=(cfl) (a)] (b) {};
                \node[r_state,fit=(cdl) (b)] (c) {};
                \node[r_state,fit=(arb) (c) ] (d) {};
            \end{tikzpicture}

        }
    \end{center}
    \caption{Иерархия языков по Хомскому}
    \label{fig:Chomsky}
\end{marginfigure}
Иерархия языков, предложенная Ноамом Хомским (Noam Chomsky), является на текущий момент классической и представлена на рисунке~\ref{fig:Chomsky}.
Она основана на сопоставлении языкам тех или иных формальных вычислителей, способных их распознать.
Например, для распознавания любого регулярного языка достаточно конечного автомата. Для контекстно-свободного~--- магазинного автомата. И так далее.

Однако, данная иерархия постепенно теряет свою актуальность, так как появляются новые классы языков, свойства которых уже не удаётся адекватным образом отобразить, используя её.

Один из вариантов иерархии языков, более полно отображающий современное состояние дел, предложен Александром Охотиным%
\sidenote{
    Иерархия и некоторые отражаемые ей свойства подробно обсуждаются в презентации Александра Охотина \enquote{Underlying principles and recurring ideas of formal grammars} (\url{https://users.math-cs.spbu.ru/~okhotin/talks/grammars_lata_talk.pdf}).
    Также, с данной презентацией рекомендуется ознакомиться чтобы представить себе состояние области в целом.
}.
Вариация предложенной Александром Охотиным иерархии\sidenote{
    Данная вариация скомпонована из версии, представленной в презентации \enquote{Underlying principles and recurring ideas of formal grammars} и версии, взятой из работы~\cite{MRYKHIN2023113829}.
} представлена на изображении~\ref{fig:hierarchyOkhotin}.
Приведённая диаграмма содержит регулярные и контекстно-свободные.
Так и подклассы, лежащие между ними.
Вместе с этим, классы, лежащие выше контекстно-свободных: многокомпонентные контекстно-свободные, булевы, конъюнктивные, их подклассы.

\begin{figure}[h!]
    \begin{center}
        \resizebox{\textwidth}{!}{
                        \begin{tikzpicture}[node distance=1.0cm,bend angle=45,auto]
                \tikzstyle{lang}=[circle,thick,draw=blue!75,fill=blue!75,minimum size=1mm]
                \node [lang] (reg) [label=below:\emph{Reg}] {};
                \node [right of = reg] (reg_dummy) {};
                \node [lang] (lllin) [right of=reg_dummy, label=below:\emph{LLLin}] {};
                \node [right of = lllin] (lllin_dummy) {};
                \node [lang] (lrlin) [right of=lllin_dummy, label=below:\emph{LRLin}] {};
                \node [right of = lrlin] (lrlin_dummy) {};
                \node [lang] (unamblin) [right of=lrlin_dummy, label=below:\emph{UnambLin}] {};
                \node [right of = unamblin] (unamblin_dummy) {};
                \node [lang] (lin) [right of=unamblin_dummy, label=below:\emph{Lin}] {};
                \node [right of = lin] (lin_dummy) {};
                \node [right of = lin_dummy] (lin_dummy_2) {};
                \node [right of = lin_dummy_2] (lin_dummy_3) {};
                \node [lang] (conj) [right of=lin_dummy_3, label=below:\emph{Conj}] {};
                \node [right of = conj] (conj_dummy) {};
                \node [lang] (bool) [right of=conj_dummy, label=right:\emph{Bool}] {};

                \node [above of = lrlin] (lrlin_up_dummy) {};
                \node [lang] (ll) [above of=lrlin_up_dummy, label=above:\emph{LL}] {};
                \node [right of = ll] (ll_dummy) {};
                \node [lang] (lr) [right of=ll_dummy, label=above:\emph{LR}] {};
                \node [right of = lr] (lr_dummy) {};
                \node [lang] (unamb) [right of=lr_dummy, label=above:\emph{Unamb}] {};
                \node [right of = unamb] (unamb_dummy) {};
                \node [lang] (ordinary) [right of=unamb_dummy, label=right:\emph{Ordinary}] {};

                \node [lang] (conjleftcontext) [above of=conj_dummy, label=right:\emph{Conj$+\lhd$}] {};

                \node [below of = lrlin] (lrlin_down_dummy) {};
                \node [lang] (vpda) [below of=lrlin_down_dummy, label=below:\emph{VPDA}] {};
                \node [below of = lin] (lin_down_dummy_1) {};
                \node [below of = lin_down_dummy_1] (lin_down_dummy_2) {};
                \node [right of = lin_down_dummy_2] (lin_right_dummy_1) {};
                \node [lang] (linconj) [right of=lin_right_dummy_1, label=below:\emph{LinConj}] {};

                \node [lang] (unambtag) [above of=unamb_dummy, label=above:\emph{UnambTAG}] {};
                \node [right of = unambtag] (unambtag_dummy) {};
                \node [lang] (tag) [right of=unambtag_dummy, label=above:\emph{TAG}] {};
                \node [right of = tag] (tag_dummy) {};
                \node [lang] (mcfl) [right of=tag_dummy, label=above:\emph{MCFL}] {};

                \node [above of = linconj] (linconj_above_dummy) {};
                \node [lang] (unambconj) [right of=linconj_above_dummy, label=below:\emph{UnambConj}] {};
                \node [right of = unambconj] (unambconj_dummy) {};
                \node [lang] (unambbool) [right of=unambconj_dummy, label=right:\emph{UnambBool}] {};

                \node [below of = linconj] (linconj_below_dummy) {};
                \node [lang] (rtca) [right of=linconj_below_dummy, label=below:\emph{RT-CA}] {};
                \node [right of = rtca] (rtca_dummy) {};
                \node [lang] (ltca) [right of=rtca_dummy, label=right:\emph{LT-CA}] {};

                \path[->]
                (reg) edge node {} (lllin)
                (lllin) edge node {} (lrlin)
                (lrlin) edge node {} (unamblin)
                (unamblin) edge node {} (lin)

                (lllin) edge node {} (ll)
                (ll) edge node {} (lr)
                (ll) edge node {} (lr)
                (lr) edge node {} (unamb)
                (unamb) edge node {} (ordinary)

                (lrlin) edge node {} (lr)
                (unamblin) edge node {} (unamb)
                (lin) edge node {} (ordinary)

                (reg) edge node {} (vpda)
                (vpda) edge node {} (lr)
                (vpda) edge node {} (linconj)

                (ordinary) edge node {} (conj)
                (conj) edge node {} (bool)
                (conj) edge node {} (conjleftcontext)

                (unamb) edge node {} (unambtag)
                (ordinary) edge node {} (tag)
                (unambtag) edge node {} (tag)
                (tag) edge node {} (mcfl)

                (lin) edge node {} (linconj)
                (linconj) edge node {} (unambconj)
                (unambconj) edge node {} (unambbool)
                (unambconj) edge node {} (conj)
                (unambbool) edge node {} (bool)
                (unamb) edge node {} (unambconj)

                (linconj) edge node {} (rtca)
                (rtca) edge node {} (ltca)
                ;
            \end{tikzpicture}

        }
    \end{center}
    \caption{Иерархия. \emph{Reg}~--- регулярные, \emph{LLLin}, \emph{LRLin}, \emph{UnambLin}, \emph{Lin}, \emph{LL}, \emph{LR}, \emph{Unamb}~--- однозначные, \emph{Ordinary}~--- обыкновенные (контекстно-свободные)
        \emph{Conj}~--- конъюнктивные, \emph{Boolean}~--- булевы, \emph{UnamnbBool}, \emph{UnambConj}, \emph{VPDA}, \emph{LinConj}, \emph{UnambTAG}, \emph{TAG}, \emph{MCFL},}
    \label{fig:hierarchyOkhotin}
\end{figure}

Для того, чтобы содержательно рассуждать про различные классы языков, необходимо иметь механизм, позволяющий чётко отделить один класс от другого.
\emph{Лемма о накачке} для соответствующего класса~--- один из классических таких механизмов.
Однако, не для всех классов языков соответствующие результаты получены.
Так, например, формулировка леммы о накачки для многокомпонентных контекстно-свободных языков в общем виде всё ещё не найдена, хотя существуют формулировки для отдельных подклассов.
Аналогично, для булевых и конъюнктивных языков всё ещё не предложены аналоги лемм о накачке.


%\section{Вопросы и задачи}
%\begin{enumerate}
%  \item !!!
%  \item !!!
%\end{enumerate}
